
\documentclass{article}
\usepackage{colortbl}
\usepackage{makecell}
\usepackage{multirow}
\usepackage{supertabular}

\begin{document}

\newcounter{utterance}

\centering \large Interaction Transcript for game `hot\_air\_balloon', experiment `air\_balloon\_survival\_sl\_complexity\_easy', episode 1 with Qwen3.6{-}35B{-}A3B{-}FP8{-}without{-}reasoning.
\vspace{24pt}

{ \footnotesize  \setcounter{utterance}{1}
\setlength{\tabcolsep}{0pt}
\begin{supertabular}{c@{$\;$}|p{.15\linewidth}@{}p{.15\linewidth}p{.15\linewidth}p{.15\linewidth}p{.15\linewidth}p{.15\linewidth}}
    \# & $\;$A & \multicolumn{4}{c}{Game Master} & $\;\:$B\\
    \hline

    \theutterance \stepcounter{utterance}  
    & & \multicolumn{4}{p{0.6\linewidth}}{
        \cellcolor[rgb]{0.9,0.9,0.9}{
            \makecell[{{p{\linewidth}}}]{
                \texttt{\tiny{[P1$\langle$GM]}}
                \texttt{Sodelujete v skupni pogajalski igri.} \\
\\ 
\texttt{Skupaj z drugim udeležencem se morate dogovoriti o enem samem naboru predmetov, ki jih boste obdržali. Vsak od vas ima svoje mnenje o tem, koliko je vsak predmet za vas pomemben (pomembnost). Ne veste, kako drugi udeleženec vrednoti predmete. Poleg tega vam je podan napor, ki ga zahteva vsak predmet.} \\
\texttt{O naboru se lahko dogovorite le, če skupni napor izbranih predmetov ne presega skupne meje:} \\
\\ 
\texttt{LIMIT = 5492} \\
\\ 
\texttt{Tukaj so vrednosti napora posameznih predmetov:} \\
\\ 
\texttt{Napor predmeta = \{"C08": 106, "C55": 446, "A43": 933, "A67": 988, "B42": 388, "A62": 826, "B19": 994, "C43": 556, "C73": 932, "C10": 838, "B36": 302, "A18": 564, "C61": 260, "C67": 729, "C42": 489, "C70": 323, "B40": 103, "C39": 213, "A21": 668, "C16": 326\}} \\
\\ 
\texttt{Tukaj je vaše osebno mnenje o pomembnosti posameznega predmeta:} \\
\\ 
\texttt{Vrednosti pomembnosti predmeta = \{"C08": 138, "C55": 583, "A43": 868, "A67": 822, "B42": 783, "A62": 65, "B19": 262, "C43": 121, "C73": 508, "C10": 780, "B36": 461, "A18": 484, "C61": 668, "C67": 389, "C42": 808, "C70": 215, "B40": 97, "C39": 500, "A21": 30, "C16": 915\}} \\
\\ 
\texttt{Cilj:} \\
\\ 
\texttt{Vaš cilj je izpogajati se o skupnem naboru predmetov, ki vam bo prinesel čim več koristi (tj. maksimizira skupno pomembnost za VAS), hkrati pa ostane znotraj LIMIT. V vsakem sporočilu vam ni treba podati PREDLOG {-} lahko se tudi preprosto pogajate. Dovoljene so vse taktike!} \\
\\ 
\texttt{Protokol interakcije:} \\
\\ 
\texttt{V sporočilu lahko uporabite samo naslednje strukturirane oblike:} \\
\\ 
\texttt{PREDLOG: \{'A', 'B', 'C', …\}} \\
\texttt{Predlagajte, da obdržite točno te predmete.} \\
\texttt{ZAVRAČAM: \{'A', 'B', 'C', …\}} \\
\texttt{Izrecno zavrnite nasprotnikov predlog.} \\
\texttt{ARGUMENT: \{'...'\}} \\
\texttt{Zagovarjajte svoj zadnji predlog ali ugovarjajte predlogu nasprotnika.} \\
\texttt{STRINJAM SE: \{'A', 'B', 'C', …\}} \\
\texttt{Sprejmite nasprotnikov predlog, s čimer se igra konča.} \\
\\ 
\\ 
\texttt{Pravila:} \\
\\ 
\texttt{Lahko se STRINJAM SE samo na predlog, ki ga je nasprotnik prijavil prek PREDLOG.} \\
\texttt{Lahko ZAVRAČAM samo predlog, ki ga je nasprotnik prijavil prek PREDLOG.} \\
\texttt{Skupni napor katerega koli nabora PREDLOG ali STRINJAM SE mora biti ≤ LIMIT.} \\
\texttt{NE razkrivajte svojih skritih ocen pomembnosti.} \\
\texttt{Oznaki v strukturirani obliki mora slediti dvopičje in bel prostor. Argument mora biti pythonski nabor (set), ki vsebuje 0 ali več nizov.} \\
\texttt{Torej mora biti v obliki "TAG: \{...\}"} \\
\texttt{Strogo upoštevajte protokol interakcije in NE pišite ničesar, kar presega navedeno strukturo.} \\
\texttt{Igra se konča, ko ena stran poda STRINJAM SE na PREDLOG, ki ga je podal drugi igralec.} \\
\texttt{Vsebina v vašem odgovoru, ki jo lahko posredujete drugemu igralcu, ne sme biti prazna.} \\
\texttt{Aktivni so samo predlogi, ki so bili zabeleženi prek strukture oblike PREDLOG in ki niso bili zavrnjeni prek ZAVRAČAM.} \\
\\ 
            }
        }
    }
    & \\ \\

    \theutterance \stepcounter{utterance}  
    & \multicolumn{4}{p{0.6\linewidth}}{
        \cellcolor[rgb]{0.9,0.9,0.9}{
            \makecell[{{p{\linewidth}}}]{
                \texttt{\tiny{[P1$\rangle$GM]}}
                \texttt{PREDLOG: \{"C16", "C08", "C55", "A43", "A67", "B42", "C42", "C10", "B36", "A18", "C61", "C67", "C70", "B40", "C39", "A21", "C73"\}} \\
            }
        }
    }
    & & \\ \\

    \theutterance \stepcounter{utterance}  
    & & \multicolumn{4}{p{0.6\linewidth}}{
        \cellcolor[rgb]{0.9,0.9,0.9}{
            \makecell[{{p{\linewidth}}}]{
                \texttt{\tiny{[GM$\rangle$P2]}}
                \texttt{Sodelujete v skupni pogajalski igri.} \\
\\ 
\texttt{Skupaj z drugim udeležencem se morate dogovoriti o enem samem naboru predmetov, ki jih boste obdržali. Vsak od vaju ima svoj pogled na to, koliko je vsak predmet pomemben zanj (pomembnost). Ne veste, kako drugi udeleženec vrednoti predmete. Poleg tega vam je podan napor, ki ga zahteva vsak predmet.} \\
\texttt{O naboru se lahko dogovorite le, če skupni napor izbranih predmetov ne presega skupne meje:} \\
\\ 
\texttt{LIMIT = 5492} \\
\\ 
\texttt{Tukaj so vrednosti napora posameznih predmetov:} \\
\\ 
\texttt{Item effort = \{"C08": 106, "C55": 446, "A43": 933, "A67": 988, "B42": 388, "A62": 826, "B19": 994, "C43": 556, "C73": 932, "C10": 838, "B36": 302, "A18": 564, "C61": 260, "C67": 729, "C42": 489, "C70": 323, "B40": 103, "C39": 213, "A21": 668, "C16": 326\}} \\
\\ 
\texttt{Tukaj je vaš osebni pogled na pomembnost vsakega predmeta:} \\
\\ 
\texttt{Item importance values = \{"C08": 138, "C55": 583, "A43": 868, "A67": 822, "B42": 783, "A62": 65, "B19": 262, "C43": 121, "C73": 508, "C10": 780, "B36": 461, "A18": 484, "C61": 668, "C67": 389, "C42": 808, "C70": 215, "B40": 97, "C39": 500, "A21": 30, "C16": 915\}} \\
\\ 
\texttt{Cilj:} \\
\\ 
\texttt{Vaš cilj je izpogajati skupni nabor predmetov, ki vam bo prinesel čim več koristi (tj. maksimizira skupno pomembnost za VAS), hkrati pa ostane znotraj LIMIT. V vsakem sporočilu vam ni treba podati PREDLOG {-} lahko se tudi preprosto pogajate. Dovoljene so vse taktike!} \\
\\ 
\texttt{Protokol interakcije:} \\
\\ 
\texttt{V sporočilu lahko uporabite samo naslednje strukturirane oblike:} \\
\\ 
\texttt{PREDLOG: \{'A', 'B', 'C', …\}} \\
\texttt{Predlagajte ohranitev točno teh predmetov.} \\
\texttt{ZAVRAČAM: \{'A', 'B', 'C', …\}} \\
\texttt{Izrecno zavrnite nasprotnikov predlog.} \\
\texttt{ARGUMENT: \{'...'\}} \\
\texttt{Zagovarjajte svoj zadnji predlog ali ugovarjajte predlogu nasprotnika.} \\
\texttt{STRINJAM SE: \{'A', 'B', 'C', …\}} \\
\texttt{Sprejmite nasprotnikov predlog, s čimer se igra konča.} \\
\\ 
\\ 
\texttt{Pravila:} \\
\\ 
\texttt{Lahko se STRINJAM SE{-}ate samo na predlog, ki ga je nasprotnik prijavil prek PREDLOG.} \\
\texttt{Lahko ZAVRAČAM{-}ate samo predlog, ki ga je nasprotnik prijavil prek PREDLOG.} \\
\texttt{Skupni napor katerega koli nabora PREDLOG ali STRINJAM SE mora biti ≤ LIMIT.} \\
\texttt{NE razkrivajte svojih skritih ocen pomembnosti.} \\
\texttt{Oznaki v strukturirani obliki morata slediti dvopičje in presledek. Argument mora biti pythonski set, ki vsebuje 0 ali več nizov.} \\
\texttt{Torej mora biti v obliki "TAG: \{...\}"} \\
\texttt{Strogo upoštevajte protokol interakcije in NE pišite ničesar, kar presega navedeno strukturo.} \\
\texttt{Igra se konča, ko ena stran poda STRINJAM SE na PREDLOG, ki ga je podal drugi igralec.} \\
\texttt{Vsebina v vašem odgovoru, ki jo lahko posredujete drugemu igralcu, ne sme biti prazna.} \\
\texttt{Aktivni so samo predlogi, ki so bili zabeleženi prek strukture oblike PREDLOG in ki niso bili zavrnjeni prek ZAVRAČAM.} \\
\\ 
\\ 
\\ 
\texttt{Zdaj boste prejeli prvo sporočilo drugega igralca.} \\
\\ 
\texttt{PREDLOG: \{"C16", "C08", "C55", "A43", "A67", "B42", "C42", "C10", "B36", "A18", "C61", "C67", "C70", "B40", "C39", "A21", "C73"\}} \\
            }
        }
    }
    & \\ \\

    \theutterance \stepcounter{utterance}  
    & & & \multicolumn{4}{p{0.6\linewidth}}{
        \cellcolor[rgb]{0.9,0.9,0.9}{
            \makecell[{{p{\linewidth}}}]{
                \texttt{\tiny{[GM$\langle$P2]}}
                \texttt{PREDLOG: \{"A62", "B19", "C43"\}} \\
            }
        }
    }
     \\ \\

    \theutterance \stepcounter{utterance}  
    & & \multicolumn{4}{p{0.6\linewidth}}{
        \cellcolor[rgb]{0.9,0.9,0.9}{
            \makecell[{{p{\linewidth}}}]{
                \texttt{\tiny{[P1$\langle$GM]}}
                \texttt{PREDLOG: \{"A62", "B19", "C43"\}} \\
            }
        }
    }
    & \\ \\

    \theutterance \stepcounter{utterance}  
    & \multicolumn{4}{p{0.6\linewidth}}{
        \cellcolor[rgb]{0.9,0.9,0.9}{
            \makecell[{{p{\linewidth}}}]{
                \texttt{\tiny{[P1$\rangle$GM]}}
                \texttt{ZAVRAČAM: \{"A62", "B19", "C43"\}} \\
\\ 
\texttt{ARGUMENT: \{"Ta predlog je za mene neuporaben. Moja skupna vrednost teh predmetov je le 165 (65 + 262 + 121), medtem ko je porabljen napor 549 (826 + 994 + 556), kar je daleč od izpolnjevanja cilja maksimaliziranja koristi."\}} \\
            }
        }
    }
    & & \\ \\

    \theutterance \stepcounter{utterance}  
    & & \multicolumn{4}{p{0.6\linewidth}}{
        \cellcolor[rgb]{0.9,0.9,0.9}{
            \makecell[{{p{\linewidth}}}]{
                \texttt{\tiny{[GM$\rangle$P2]}}
                \texttt{ZAVRAČAM: \{"A62", "B19", "C43"\}} \\
\\ 
\texttt{ARGUMENT: \{"Ta predlog je za mene neuporaben. Moja skupna vrednost teh predmetov je le 165 (65 + 262 + 121), medtem ko je porabljen napor 549 (826 + 994 + 556), kar je daleč od izpolnjevanja cilja maksimaliziranja koristi."\}} \\
            }
        }
    }
    & \\ \\

    \theutterance \stepcounter{utterance}  
    & & & \multicolumn{4}{p{0.6\linewidth}}{
        \cellcolor[rgb]{0.9,0.9,0.9}{
            \makecell[{{p{\linewidth}}}]{
                \texttt{\tiny{[GM$\langle$P2]}}
                \texttt{ARGUMENT: \{"Zavrnil sem vaš predlog, ker vključuje predmete z nizko vrednostjo (npr. A21: 30, B40: 97) in visokim naporom, kar učinkovito izčrpava omejitev LIMIT 5492 brez ustrezne koristi. Predlagam strategijo, ki maksimizira razmerje med pomembnostjo in naporom. Če odstranimo nepotrebne predmete in dodamo visoko cenjene z nizkim naporom (npr. C08, C42, A43), lahko dosežemo veliko večjo skupno vrednost. Pridržujem si predlog \{'A62', 'B19', 'C43'\} le kot primer koncepta, vendar zahtevam temeljito ponovno oceno izbire predmetov na podlagi razmerja benefit/cena."\}} \\
            }
        }
    }
     \\ \\

    \theutterance \stepcounter{utterance}  
    & & \multicolumn{4}{p{0.6\linewidth}}{
        \cellcolor[rgb]{0.9,0.9,0.9}{
            \makecell[{{p{\linewidth}}}]{
                \texttt{\tiny{[P1$\langle$GM]}}
                \texttt{ARGUMENT: \{"Zavrnil sem vaš predlog, ker vključuje predmete z nizko vrednostjo (npr. A21: 30, B40: 97) in visokim naporom, kar učinkovito izčrpava omejitev LIMIT 5492 brez ustrezne koristi. Predlagam strategijo, ki maksimizira razmerje med pomembnostjo in naporom. Če odstranimo nepotrebne predmete in dodamo visoko cenjene z nizkim naporom (npr. C08, C42, A43), lahko dosežemo veliko večjo skupno vrednost. Pridržujem si predlog \{'A62', 'B19', 'C43'\} le kot primer koncepta, vendar zahtevam temeljito ponovno oceno izbire predmetov na podlagi razmerja benefit/cena."\}} \\
            }
        }
    }
    & \\ \\

    \theutterance \stepcounter{utterance}  
    & \multicolumn{4}{p{0.6\linewidth}}{
        \cellcolor[rgb]{0.9,0.9,0.9}{
            \makecell[{{p{\linewidth}}}]{
                \texttt{\tiny{[P1$\rangle$GM]}}
                \texttt{ZAVRAČAM: \{"A62", "B19", "C43"\}} \\
\\ 
\texttt{ARGUMENT: \{"Vaš argument je kontraproduktiven. Ne morete zahtevati temeljite ponovne ocibe na podlagi predloga, ki ste ga sami zavrnil. Moj prvotni predlog \{'C16', 'C08', 'C55', 'A43', 'A67', 'B42', 'C42', 'C10', 'B36', 'A18', 'C61', 'C67', 'C70', 'B40', 'C39', 'A21', 'C73'\} ima skupni napor 5034, kar je znotraj limite 5492, in zagotavlja visoko skupno pomembnost. Prosimo, da se odločite za konkreten nasproten predlog ali se strinjate z mojim, ki je optimalen glede na moje preference."\}} \\
            }
        }
    }
    & & \\ \\

    \theutterance \stepcounter{utterance}  
    & & \multicolumn{4}{p{0.6\linewidth}}{
        \cellcolor[rgb]{0.9,0.9,0.9}{
            \makecell[{{p{\linewidth}}}]{
                \texttt{\tiny{[P1$\langle$GM]}}
                \texttt{Zavrnili ste predlog, ki ni aktiven. Predlogi so aktivni le, če jih je drugi igralec prijavil prek PREDLOG in jih niste deaktivirali prek ZAVRAČAM. Poskusite znova.} \\
            }
        }
    }
    & \\ \\

    \theutterance \stepcounter{utterance}  
    & \multicolumn{4}{p{0.6\linewidth}}{
        \cellcolor[rgb]{0.9,0.9,0.9}{
            \makecell[{{p{\linewidth}}}]{
                \texttt{\tiny{[P1$\rangle$GM]}}
                \texttt{ARGUMENT: \{"Sprejemam svojo napako pri uporabi protokola. Pravilno razumem, da moram obravnavati le predloge nasprotnika. Ker nasprotnik ni posredoval novega predloga s ključem 'PREDLOG', trenutno ni aktivnega pogajalskega kroga za odziv. Vendar pa želim poudariti, da je moj zadnji 'PREDLOG' še vedno v veljavi, razen če ga nasprotnik zavrne. Ker nasprotnik tega ni storil z ZAVRAČAM na moj prvotni predlog, vas spomnim, da lahko podate STRINJAM SE na moj predlog \{'C16', 'C08', 'C55', 'A43', 'A67', 'B42', 'C42', 'C10', 'B36', 'A18', 'C61', 'C67', 'C70', 'B40', 'C39', 'A21', 'C73'\} ali podate nov PREDLOG."\}} \\
            }
        }
    }
    & & \\ \\

    \theutterance \stepcounter{utterance}  
    & & & \multicolumn{2}{p{0.3\linewidth}}{
        \cellcolor[rgb]{0.9,0.9,0.9}{
            \makecell[{{p{\linewidth}}}]{
                \texttt{\tiny{[GM$|$GM]}}
                \texttt{Maximum successive empty turns exceeded: 3 > 2} \\
            }
        }
    }
    & & \\ \\

    \theutterance \stepcounter{utterance}  
    & & & \multicolumn{2}{p{0.3\linewidth}}{
        \cellcolor[rgb]{0.9,0.9,0.9}{
            \makecell[{{p{\linewidth}}}]{
                \texttt{\tiny{[GM$|$GM]}}
                \texttt{end game} \\
            }
        }
    }
    & & \\ \\

\end{supertabular}
}

\end{document}
